\documentclass{standalone}
\usepackage{circuitikz}
\usetikzlibrary{calc}
\definecolor{circuitnotemark}{HTML}{FE00FE}
\begin{document}
\begin{circuitikz}[american, line width=0.8pt]
\ctikzset{bipoles/length=1.2cm}
\ctikzset{grounds/scale=1.36}
\foreach \x in {0,2,4,6,8,10} {\foreach \y in {0,-2,-4,-6,-8,-10,-12} {\fill[gray, opacity=0.35] (\x,\y) circle (1.1pt);}}
\node[gray, font=\scriptsize] at (0,0.84) {1};
\node[gray, font=\scriptsize] at (2,0.84) {2};
\node[gray, font=\scriptsize] at (4,0.84) {3};
\node[gray, font=\scriptsize] at (6,0.84) {4};
\node[gray, font=\scriptsize] at (8,0.84) {5};
\node[gray, font=\scriptsize] at (10,0.84) {6};
\node[gray, font=\scriptsize, anchor=east] at (-0.84,0) {1};
\node[gray, font=\scriptsize, anchor=east] at (-0.84,-2) {2};
\node[gray, font=\scriptsize, anchor=east] at (-0.84,-4) {3};
\node[gray, font=\scriptsize, anchor=east] at (-0.84,-6) {4};
\node[gray, font=\scriptsize, anchor=east] at (-0.84,-8) {5};
\node[gray, font=\scriptsize, anchor=east] at (-0.84,-10) {6};
\node[gray, font=\scriptsize, anchor=east] at (-0.84,-12) {7};
\coordinate (x3y1) at (4,0);
\coordinate (x3y3) at (4,-4);
\coordinate (x3y5) at (4,-8);
\coordinate (x3y7) at (4,-12);
\coordinate (x6y3) at (10,-4);
\node[vcc] at (x3y1) {VCC}; % line 3
\draw (x3y1) to[R, l_=$R_{1}$, a^=$10\,\mathrm{k}\Omega$] (x3y3); % line 4
\node[npn, photo, nobase] (part-Q1) at (x3y5) {}; % line 5
\node[anchor=west] at (part-Q1.east) {$Q_{1}$}; % line 5
\node[ground] at (x3y7) {}; % line 6
\draw (x6y3) node[ocirc]{} node[above left]{OUT}; % line 7
\draw (x3y3) -- (part-Q1.collector); % line 9
\draw (part-Q1.emitter) -- (x3y7); % line 10
\draw (x3y3) -- (x6y3); % line 11
\node[circ] at (x3y3) {};
\node[anchor=south west, inner sep=2pt, yshift=4pt, circuitnotemark, font=\LARGE\bfseries] (circuittitle) at (current bounding box.north west) {X};
\path (circuittitle.south west) rectangle ++(10.058,0.853);
\node[anchor=north east, inner sep=2pt, circuitnotemark, font=\large] (circuitstamp) at ([yshift=-0.593cm]current bounding box.south east) {X};
\path (circuitstamp.north east) rectangle ++(-5.08,-0.593);
\end{circuitikz}
\end{document}
